% ===========================================================================
%  Capítulo 5 — Detalles de implementación
% ===========================================================================
\chapter{Detalles de implementaci\'on}
\label{cap:implementacion}

Este cap\'{\i}tulo describe los aspectos internos del sistema: la organizaci\'on
de los archivos con sus responsabilidades, los macros que implementan el DSL,
el patr\'on de extensibilidad del backend y la validaci\'on del sistema sobre
los tres casos de uso presentados en el cap\'{\i}tulo~\ref{cap:dominio}.

% ---------------------------------------------------------------------------
\section{Organizaci\'on en cuatro archivos}
\label{sec:organizacion}
% ---------------------------------------------------------------------------

El sistema est\'a organizado en cuatro archivos Lisp con responsabilidades
estrictamente separadas.  La dependencia entre ellos es un\'{\i}voca:
cada archivo carga \'unicamente al que le precede.

\begin{table}[h!]
\centering
\caption{Archivos del sistema y sus responsabilidades.}
\label{tab:archivos}
\renewcommand{\arraystretch}{1.3}
\begin{tabular}{p{4.8cm} p{9.0cm}}
\toprule
\textbf{Archivo} & \textbf{Responsabilidad} \\
\midrule
\file{codigo-tesis.lisp}
  & Utilidades generales: \dsl{symb}, \dsl{mkstr}, \dsl{keywd},
    \dsl{ppme} y la macro \dsl{defclass*}.  No conoce ning\'un concepto
    del dominio de hojas de c\'alculo. \\
\midrule
\file{ast-def.lisp}
  & Definici\'on de todas las clases del AST con \dsl{defclass*}.
    Backend-agn\'ostico: no contiene ninguna letra de columna Excel,
    cadena de f\'ormula ni referencia a ninguna biblioteca de salida. \\
\midrule
\file{code-generation-utils.lisp} \newline
\file{generate-code-direct.lisp}
  & Backend Excel/Python.  \'Unico lugar del sistema donde viven letras
    de columna, cadenas de f\'ormula, estructuras openpyxl y conocimiento
    de la disposici\'on concreta del workbook.  El primero contiene
    las utilidades y variables din\'amicas; el segundo, los m\'etodos
    \dsl{compile-excel-formula}. \\
\midrule
\file{dsl-directo.lisp}
  & Macros del DSL: \dsl{def-table}, \dsl{tabla}, \dsl{hoja}, \dsl{hoja-v},
    \dsl{libro}, \dsl{exists}, \dsl{conditional-rendering}, \dsl{nav},
    \dsl{col}, \dsl{tcol}, \dsl{trange}, \dsl{countif}, \dsl{counta},
    \dsl{sum-range}, \dsl{str}, \dsl{concat}, \dsl{time-add}, y dem\'as.
    No contiene l\'ogica de generaci\'on de c\'odigo. \\
\bottomrule
\end{tabular}
\end{table}

Esta separaci\'on tiene consecuencias pr\'acticas importantes: es posible
cargar \file{ast-def.lisp} y construir un AST sin haber cargado el backend,
lo que facilita las pruebas de la sintaxis y la sem\'antica del DSL de forma
independiente.

% ---------------------------------------------------------------------------
\section{Macros del DSL}
\label{sec:macros-dsl}
% ---------------------------------------------------------------------------

Los macros del DSL se clasifican en tres grupos seg\'un su papel en la
jerarqu\'{\i}a del AST.

% ............................................................................
\subsection{Macros de expresi\'on}
\label{ssec:macros-expresion}
% ............................................................................

Los macros de expresi\'on construyen nodos de la familia de expresiones de
celda.  Todos son equivalencias directas: reciben los argumentos del
programador y los pasan al constructor CLOS correspondiente, haciendo el
quoting adecuado para los s\'{\i}mbolos del dominio.

\begin{lstlisting}[language=CommonLisp,
  caption={Macros de expresi\'on representativos.}]
;; (col nombre) -> xl-expr-column-ref con nombre quoteado
(defmacro col (name &optional context)
  (if context
      `(xl-expr-column-ref :name ',name :context ,context)
      `(xl-expr-column-ref :name ',name :context nil)))

;; (trange tabla desde) / (trange tabla desde hasta)
;; -> xl-table-range con los dos extremos como xl-expr-table-col-ref
(defmacro trange (table-id from &optional to)
  `(xl-table-range
     :table-id ',table-id
     :from-col (xl-expr-table-col-ref :table-id ',table-id
                                       :name ',from :context nil)
     :to-col   ,(if to
                    `(xl-expr-table-col-ref :table-id ',table-id
                                             :name ',to :context nil)
                    nil)))

;; (exists (var :rows-of tabla) :self-in (...) :matching (...))
;; -> xl-expr-exists con domain cross-table
(defmacro exists ((var &rest domain-spec)
                  &key matching any-overlap self-in)
  (let ((domain-type (first domain-spec)))
    (ecase domain-type
      (:all-rows
       `(xl-expr-exists :bind-var     ',var
                         :domain       (xl-domain-rows :table-id nil)
                         :match-keys   ',matching
                         :overlap-cols ',any-overlap
                         :self-in-cols  nil))
      (:rows-of
       (let ((table-id (second domain-spec)))
         `(xl-expr-exists :bind-var     ',var
                           :domain       (xl-domain-rows
                                           :table-id ',table-id)
                           :match-keys   ',matching
                           :overlap-cols ',any-overlap
                           :self-in-cols  ',self-in))))))
\end{lstlisting}

La regla de quoting de los macros es determinista: un s\'{\i}mbolo que es
nombre de columna o tabla en el dominio se quotea (\dsl{',sym}); un
s\'{\i}mbolo que es una variable Lisp en scope no se quotea (\dsl{,sym}).
Esta distinci\'on es expl\'{\i}cita en todos los macros.

% ............................................................................
\subsection{Macros de estructura}
\label{ssec:macros-estructura}
% ............................................................................

Los macros \dsl{def-table}, \dsl{hoja}, \dsl{hoja-v} y \dsl{libro}
construyen los nodos de la jerarqu\'{\i}a del workbook.

\dsl{def-table} es el m\'as complejo: genera una funci\'on que construye
un \dsl{xl-table} con sus slots rellenos.  Los \dsl{:inst-params} declarados
se a\~naden a la firma de esa funci\'on, lo que permite instanciar la misma
plantilla de tabla con valores de par\'ametro distintos:

\begin{lstlisting}[language=CommonLisp,
  caption={Expansi\'on simplificada de def-table.}]
(def-table aulas-dia-table
  ((dia "") (aula1 "") ...)
  :inst-params (dia)
  :computed ((aula1 (collect-over *grupos-all* (g) ...))))

;; Genera la funcion:
(defun aulas-dia-table (&key data params dia)
  (xl-table :id 'aulas-dia-table
            :col-names '(dia aula1 ...)
            :computed (list (cons 'aula1 (collect-over ...)))
            ...))

;; Que luego se instancia con:
(tabla aulas-dia-table :dia lun :data *datos-aulas-lunes*)
;; -> (aulas-dia-table :data *datos-aulas-lunes* :dia 'lun)
\end{lstlisting}

% ............................................................................
\subsection{El macro \texttt{nav} y el sistema de posicionamiento}
\label{ssec:macro-nav}
% ............................................................................

El macro \dsl{nav} convierte una secuencia plana de pares direcci\'on--cantidad
en una lista de cons-celdas:

\begin{lstlisting}[language=CommonLisp,
  caption={Implementaci\'on del macro nav.}]
(defmacro nav (anchor &rest steps)
  (let ((step-pairs
         (loop for (dir n) on steps by #'cddr
               collect `(cons ',dir ,n))))
    `(xl-nav :anchor ,anchor
             :steps (list ,@step-pairs))))

;; Ejemplo: (nav (ultima-fila t c) abajo 1 derecha 2)
;; Expande a:
;; (xl-nav :anchor (xl-anchor :table-id 't :col-name 'c
;;                              :anchor-type :ultima-fila)
;;          :steps (list (cons 'abajo 1) (cons 'derecha 2)))
\end{lstlisting}

El backend resuelve los pasos en tiempo de generaci\'on: dado el n\'umero de
fila de la \'ultima celda de datos, suma o resta las cantidades seg\'un la
direcci\'on y produce la coordenada Excel absoluta.

% ---------------------------------------------------------------------------
\section{Extensibilidad del sistema}
\label{sec:extensibilidad}
% ---------------------------------------------------------------------------

La extensibilidad es una consecuencia directa del uso de CLOS y de la
separaci\'on en archivos.  Considrese el ejemplo de a\~nadir soporte para
la funci\'on \texttt{AVERAGEIF} de Excel, que actualmente no est\'a en el
lenguaje.

\textbf{Paso 1 --- Nuevo nodo AST} (en \file{ast-def.lisp}):
\begin{lstlisting}[language=CommonLisp]
(defclass* xl-expr-averageif () (average-range criteria))
\end{lstlisting}

\textbf{Paso 2 --- Nuevo macro} (en \file{dsl-directo.lisp}):
\begin{lstlisting}[language=CommonLisp]
(defmacro averageif (range-expr criteria)
  `(xl-expr-averageif :average-range ,range-expr
                       :criteria ,criteria))
\end{lstlisting}

\textbf{Paso 3 --- Nuevo m\'etodo} (en \file{generate-code-direct.lisp}):
\begin{lstlisting}[language=CommonLisp]
(defmethod compile-excel-formula
    ((e clase-xl-expr-averageif) col-map data-names
     row-num first-row last-row)
  (let ((range-str (compile-excel-formula
                     (average-range e) col-map data-names
                     row-num first-row last-row))
        (criteria  (compile-excel-formula
                     (criteria e) col-map data-names
                     row-num first-row last-row)))
    (format nil "AVERAGEIF(~a,~a)" range-str criteria)))
\end{lstlisting}

Ning\'un c\'odigo existente necesita modificarse.  El macro \dsl{defclass*}
garantiza que la nueva clase tiene exactamente la misma interfaz que las dem\'as.

% ---------------------------------------------------------------------------
\section{Validaci\'on sobre los casos de uso}
\label{sec:validacion}
% ---------------------------------------------------------------------------

Los tres casos de uso descritos en el cap\'{\i}tulo~\ref{cap:dominio} se
validaron ejecutando el DSL y verificando manualmente el archivo
\xlsx{} generado.

% ............................................................................
\subsection{Horario de Defensas de Tesis}
\label{ssec:val-defensas}
% ............................................................................

El programa DSL del listado~\ref{lst:tutorial} genera \texttt{Horario\_Tesis.xlsx}
con la estructura esperada.  La verificaci\'on confirm\'o:

\begin{itemize}
  \item La columna \texttt{Defensas} de \dsl{profesores-table} muestra los
    valores correctos: Piad~=~2, Garc\'{\i}a~=~2, L\'opez~=~3, Torres~=~3,
    Ruiz~=~2, Soto~=~2.
  \item Los profesores L\'opez y Torres, que participan en las dos defensas
    del slot \texttt{Lunes~10:00}, aparecen resaltados en rojo.
  \item Los totales en la fila inferior a las tablas son correctos:
    ``Total defensas: 3'' y ``Total profesores: 6'' con las sumas esperadas.
  \item La \texttt{FormulaRule} persiste en el archivo: al cambiar
    manualmente la hora de la tercera defensa de \texttt{10:00} a
    \texttt{14:00}, el resaltado de L\'opez y Torres desaparece
    autom\'aticamente en Excel.
\end{itemize}

% ............................................................................
\subsection{Horario Docente de la Facultad}
\label{ssec:val-facultad}
% ............................................................................

El programa \file{ast-facultad.lisp} genera \texttt{Horario\_Facultad.xlsx}
con diecisiete hojas de grupo (Ciencia de la Computaci\'on en cuatro a\~nos,
Matem\'atica en cuatro a\~nos y Ciencia de Datos en cuatro a\~nos) m\'as una
hoja de ocupaci\'on de aulas.

La verificaci\'on confirm\'o:

\begin{itemize}
  \item Las columnas \dsl{asignadas} y \dsl{faltan} de \dsl{stats-table}
    calculan correctamente para cada grupo a partir del contenido de
    \dsl{turno-table}.
  \item La hoja de aulas muestra en cada celda los grupos que ocupan cada
    aula en cada turno, con la lista de grupos separada por comas producida
    por las f\'ormulas \texttt{SUBSTITUTE(TRIM(\ldots))}.
  \item Las celdas compuestas de \dsl{turno-table} (\dsl{cell-height~=~2})
    se generan correctamente: la subfila superior muestra la asignatura y
    la inferior el aula.
\end{itemize}

Este caso ejercita, adem\'as, la caracter\'{\i}stica de \dsl{inst-params}:
la tabla \dsl{aulas-dia-table} se define una sola vez y se instancia una
vez por d\'{\i}a de la semana, con \dsl{:dia lun}, \dsl{:dia mar}, etc.

% ............................................................................
\subsection{Discusi\'on comparativa}
\label{ssec:val-comparacion}
% ............................................................................

Para ilustrar el beneficio del DSL, se estima la cantidad de c\'odigo
necesaria para implementar el caso de defensas de forma imperativa directa
en Python con openpyxl, en comparaci\'on con el programa DSL.

La versi\'on imperativa directa requiere calcular manualmente las letras de
columna, construir las cadenas de f\'ormula COUNTIF con interpolaci\'on de
texto, y ensamblar la f\'ormula SUMPRODUCT con referencias absolutas calculadas
a partir del n\'umero de filas.  Una implementaci\'on representativa ocupa
aproximadamente 80--100 l\'{\i}neas de Python.

El programa DSL equivalente (listado~\ref{lst:tutorial}) ocupa 40 l\'{\i}neas,
de las cuales la mayor parte son la declaraci\'on de datos
(\dsl{def-table}, \dsl{libro}) y la expresi\'on de intenci\'on
(\dsl{exists}, \dsl{conditional-rendering}).  Ning\'una l\'{\i}nea
contiene una letra de columna, un \'{\i}ndice de fila o una cadena de f\'ormula.

M\'as importante que el n\'umero de l\'{\i}neas es que el programa DSL es
resistente a cambios estructurales: a\~nadir una columna a \dsl{defensas-table}
s\'olo requiere a\~nadir un par \dsl{(nombre "Cabecera")} a la lista de columnas;
el compilador recalcula autom\'aticamente todas las letras y rangos.  En la
versi\'on imperativa, el mismo cambio requiere actualizar manualmente todas
las referencias a columnas afectadas.
